\chapter{Cross-level solution consistency (menu 7)}
\label{ch:menu7}
\menulabel{7}

\section{Purpose}

When the same molecule has several candidate solutions (different 5f
occupations, different SCF branches) each computed at a cheap level (say HF)
and an expensive level (say CCSD(T)), check whether ``pick the best solution by
the cheap level'' is safe. The classic f-block failure mode: the ordering of
solutions changes with the method.

\section{What you need}

A records JSON: one entry per candidate solution with a label and one energy
per level. The top level may be a list or an object with a \code{records} list;
extra fields such as \code{occupation} are kept for readability and otherwise
ignored. The schema is in Chapter~\ref{ch:formats}.

\section{Walkthrough}

\lstinputlisting[style=fbkcode, firstline=16, lastline=25,
  title={The menu-7 example session on the bundled PuCl\textsubscript{3} table}]
  {examples/expected/m7_crosslevel.txt}

\section{What you get}

Warnings, each with its evidence: when the two levels disagree about the best
solution, and when a solution inside the cheap level's top $N$ drops out of the
expensive level's top $N$. For the bundled example (a literature table of 21
PuCl$_3$ 5f occupations with HF and CCSD(T) energies, taken from the SI of the
source paper), both fire:

\begin{itemize}
  \item HF favours solution No.~3, CCSD(T) favours No.~1 -- picking by HF
    leaves the pick 0.06 kcal/mol behind at the CCSD(T) level (small, but the
    ordering is what is wrong);
  \item the HF second-best solution (No.~10) drops to position 7 at the
    CCSD(T) level (about 1.76 kcal/mol behind) -- a much larger move, and the
    qualitative lesson.
\end{itemize}

\section{How to read the numbers}

\begin{readbox}
The absolute energy differences in the example are small; the transferable
lesson is the ordering instability. Two consequences: re-rank the candidates at
the expensive level before choosing one, and re-check the identity of a
solution that moves (a dropped solution may be a different electronic state
re-labelled by the occupation table).
\end{readbox}

\section{Worked example}

The bundled file \file{fixtures/literature/pucl3\_s18.json} is the worked
example: prepare your own records in the same shape from two levels you
already ran, run \menu{7} on it, and compare the verdicts with your chemical
expectation. If the warnings fire on your data, the cheap-level shortlist is
not a safe selector.

\section{Caveats, refusals and related sections}

\begin{refusalbox}
\begin{itemize}
  \item The check works on energies you supply; it does not compute anything
    and does not decide which solution is physically right.
  \item The top-$N$ depth is a program constant; the tool states which $N$ it
    used in the output when the rule fires.
  \item Malformed records (missing labels, non-numeric energies, fewer than
    two levels) are refused with a next-step message.
\end{itemize}
\end{refusalbox}

Related: \menu{9} for SCF-side multiple-solution problems, and the A1 section
of \menu{1} -- ``no f-composition orbital'' is the same phenomenon seen from
one output file.
